\section{Theoretische Grundlagen}

Zunächst sollen also die theoretischen Grundlagen der Triangulation dargelegt werden. Zum Verständnis der methodischen Triangulation und zur späteren Verwendung in der Analyse der ausgewählten Fallstudie, soll auf die quantitative und qualitative Forschungsmethodik, sowie ihre Kombination im Mixed-Methods-Ansatz eingegangen werden.

\subsection{Quantitative Forschungsmethodik}

Die Quantitative Forschungsmethodik zielt auf einen Erkenntnisgewinn aus einer möglichst breiten, repräsentativen Datenbasis, die durch standardisierte, objektive Verfahren erhoben wurde. Besonders häufig werden dabei Umfragen, Experimente oder Beobachtungen eingesetzt. Die Daten werden in der Regel numerisch erfasst und statistisch ausgewertet, um Zusammenhänge, Unterschiede oder Trends zu identifizieren. Ein Vorteil der quantitativen Forschungsmethodik liegt in der Möglichkeit, Ergebnisse zu generalisieren und Hypothesen zu testen. Allerdings kann sie auch Einschränkungen haben, insbesondere wenn es darum geht, komplexe Phänomene oder subjektive Erfahrungen zu erfassen.

\subsection{Qualitative Forschungsmethodik}

Dem gegenüber steht die Qualitative Forschungsmethodik, die darauf abzielt, ein tieferes Verständnis von sozialen Phänomenen zu erlangen, indem sie sich auf die subjektiven Erfahrungen und Perspektiven der Teilnehmer konzentriert. Qualitative Methoden umfassen beispielsweise Interviews, Fokusgruppen oder Grounded Theory. Die Daten werden in der Regel in Form von Texten, Bildern oder Audioaufnahmen erfasst und durch interpretative Verfahren analysiert. Ein Vorteil der qualitativen Forschungsmethodik liegt in ihrer Fähigkeit, komplexe soziale Prozesse und Bedeutungen zu erfassen. Allerdings kann sie auch Einschränkungen haben, insbesondere wenn es darum geht, Ergebnisse zu generalisieren oder objektiv zu messen.

\subsection{Mixed-Methods-Design}

Werden beide Forschungsparadigma im Rahmen einer Studie kombiniert, so spricht man von einem Mixed-Methods-Design. Dabei können je nach Reihenfolge bzw. Beziehung der Methodiken unterschiedliche Zwecke verfolgt werden. Eine qualitative Analyse kann verwendet werden, um die Ergebnisse aus der quantitativen Erhebung zu vertiefen und besser zu verstehen. Dabei spricht man nach der Einteilung von Creswell und Plano Clark (2018) TODO von einem sequenziell erklärendem Design. Beginnt man hingegen mit quantitativer Analyse können ihre Ergebnisse der explorativen Theorienbildung dienen, die anschließend durch eine quantitative Methode überprüft werden kann. Diese Reihenfolge bei der mit der qualitativen Analyse begonnen wird, bezeichnet man als sequenzielles exploratives Design. Neben diesen beiden Formen des sequenziellen Designs existieren ebenfalls parallele Designs, bei denen beide Forschungsparadigmen in etwa zur gleichen Zeit angewendet werden. (Jens Winkel,... TODO) Grundsätzlich benötigt die Durchführung mehrerer Forschungsmethoden einen größeren Aufwand, kann dafür aber die Stärken beider Paradigmen kombinieren und sowohl eine inhaltliche Tiefe als auch eine statistische Breite erreichen.

\subsection{Methodische Triangulation}

In Anlehnung an die geodätische Praxis geht es nach Norman Denzin in der Triangulation, um einen Erkenntnisgewinn durch die Kombination von Messungen aus mehreren unterschiedlichen Perspektiven. Die methodische Triangulation wendet dieses Konzept auf die Forschungsmethodik an und kombiniert unterschiedliche Methoden. Neben der Methodentriangulation können auch Daten (Datentriangulation), Theorien (Theorietriangulation) oder Standpunkte verschiedener Forscher (Investigator-Triangulation) kombiniert werden.

\subsection{Bedeutung von Triangulation für Validität und Zuverlässigkeit}

Die Triangulation kann eingesetzt werrden, um die Validität und Zuverlässigkeit von Forschungsergebnissen zu steigern. Die Validität wird allgemein unterteilt in Konstruktvalidität, interne Validität und externe Validität. Die Konstruktvalidität beschreibt, ob die Messungen tatsächlich die theoretischen Konstrukte messen können, die sie messen sollen. Dieser Fehler kann durch die Kombination von Methoden reduziert werden, da unterschiedliche Methoden unterschiedliche Aspekte eines Konstrukts erfassen können. Die interne Validität beschreibt, ob die Ergebnisse tatsächlich auf die untersuchten Variablen zurückzuführen sind oder ob andere Faktoren eine Rolle spielen. Auch hier kann eine weitere Methode die möglicherweise falschen Rückschlüsse korrigieren. Die externe Validität beschreibt, ob die Ergebnisse auf andere Situationen oder Populationen generalisiert werden können. Dies kann besonders durch die Kombination mit quantitativer Methodik erreicht werden.

Die Zuverlässigkeit (Reliabilität) besteht aus Reproduzierbarkeit und Replizierbarkeit. Die Reproduzierbarkeit wird allgemein als ein Vorteil quantitativer Methotik anerkannt. Daher bewirkt methodische Triangulation hier nur eine Erhöhung der Reproduzierbarkeit, wenn eine qualitative Analyse durch eine quantitative Messung erweitert wird. Daten- und Investigator-Triangulation können Reproduzierbarkeit direkt steigern, da es bei Reproduzierbarkeit schlielich genau darum dreht, dass auch andere Daten und andere Forscher zu gleichen Ergebnissen kommen. Triangualtion egal welcher Form kann die Replizierbarkeit steigern, da die Ergebnisse dabei aus mehreren Quellen kommen, statt sich lediglich auf ein Ergebnis einer Person mit einer Methode und einer Datenbasis bezieht.

% TODO: Quelle, Korrektheit prüfen, ist das schon zu viel eigene Meinung hier?

% TODO: ist der Abschnitt zu Validität und Zuverlässigkeit nicht schon etwas vorwegnehmend? -> relativieren: „Triangulation wird in der Literatur mit dem Ziel eingesetzt, Validität und Zuverlässigkeit zu erhöhen – ob und wie das gelingt, ist jedoch kontextabhängig."

% TODO: Morse Design Notation ergänzen